\begin{abstract}
\noindent Diabetic retinopathy is a leading cause of preventable vision loss, and automated screening usually ignores its earliest, microaneurysm-only stage. This study aimed to build a lesion-aware deep learning system for detecting any diabetic retinopathy, with attention to the mild stage, and to audit the public data used to evaluate such systems. Four public datasets were audited for duplicates, conflicting grades, image quality and lesion size. A quality gate and a classifier with learnable lesion queries and a class-weighted focal loss were built on a partially fine-tuned EfficientNet-B0, developed on 11,998 EyePACS images, and compared with a global-pooling baseline on an internal test set and on three external sets, with bootstrap confidence intervals. The audit found 148 duplicate pairs in APTOS, 43 with different grades, and a duplicate crossing the official IDRiD train/test split; half of all microaneurysms were narrower than one pixel at a 224-pixel input, and 61% of proliferative images were labelled ungradable. The proposed system reached sensitivities of 0.91 to 0.98 and detected 84% to 91% of mild cases, significantly more than the baseline on two sets, but its specificity was 0.10 to 0.21 lower and its area under the curve slightly lower. Neither the attention maps nor Grad-CAM located lesions better than chance. Early detection is limited more by input resolution and data quality than by the classifier head, and quality rejection must trigger referral rather than a negative result.

\medskip\noindent\textbf{Keywords:} diabetic retinopathy; early detection; deep learning; external validation; dataset audit; explainability
\end{abstract}
